% sections/oracle.tex
\section{C1: How well full-length rehearsal predicts the leaderboard}
\label{sec:oracle}

\subsection{Protocol}

The rule we followed for every rehearsed upload from 26~September on (all except the four official tests
\upl{K65a}--\upl{K65d}) is short:
\begin{quote}
  Run the \emph{exact upload file}, at the upload seed, on your own Trainium instance, for the full charged budget,
  starting \emph{cold}. Score the trained weights on the first \val{eval-tokens} tokens of the public validation
  data. The official score will be that number plus an offset.
\end{quote}
Each shuffle salt is a code constant, so every salt is its own file with its own hash, and its rehearsal is a
rehearsal of that file. Among the rehearsed uploads, two depart from the rule. The rehearsal of \upl{K70a} ran on
chip~E, which is slow, and was normalised to chip-C speed before use (\cref{sec:k70a}). \upl{K54b} was rehearsed
before its comments were stripped to bring it under the size limit; the uploaded file, \upl{K54b\_min}, parses to the
\val{k54b-ast}\docnum{} as the rehearsed one, so it runs the same program. For upload $i$ with rehearsal
score $r_i$ on chip $c(i)$ and official score $y_i$ we write
\begin{equation}
  \label{eq:offset}
  y_i = r_i + \mu + \gamma_{c(i)} + \varepsilon_i, \qquad \varepsilon_i \sim \mathcal{N}(0, \sigma_{\text{off}}^2),
\end{equation}
where $\mu$ is the part of the offset that belongs to the evaluation text and $\gamma_c$ a per-chip part. During
the campaign we used the special case $\gamma\equiv0$: estimate $\hat\mu$ and $s$ from $n$ earlier uploads and
predict $y_{\text{new}}$ with the interval $r_{\text{new}} + \hat\mu \pm t_{n-1,0.975}\, s\sqrt{1 + 1/n}$
(\cref{app:derivations}). Apart from \upl{K65a}--\upl{K65d}, a file was uploaded only after its cold rehearsal
had exited cleanly, with no graph breaks, charged time inside the budget and the file under the size limit. The
repository records written projections for three uploads (\cref{sec:projections}).

\paragraph{Why it should work, and why not exactly.} If training is deterministic (\cref{sec:determinism}), the
official run follows the rehearsal's trajectory for as many steps as it gets. The text then contributes a
near-constant: the first 2M public tokens are easier than the first 20M by \val{gap-mean} (SD \val{gap-sd} over
\val{gap-n} runs; \cref{sec:contest}), which is the size of the offset. But the official host need not run at the rehearsal chip's
speed. A host that is $x\%$ faster or slower gives $x\%$ more or fewer steps, worth $\rho x$ at the price of
\cref{sec:price}. That is the term $\gamma_c$: it is a property of the \emph{rehearsal chip}, relative to the
official host, not of the text.

\subsection{What the offset is made of}

\Cref{tab:offset-periods} (\cref{app:oracle}) gives the offset by period. Since \upl{K50} the official score has been the rehearsal plus
\val{off-mean-since} (SD \val{off-sd-since}, range \val{off-min-since} to \val{off-max-since},
$n=\val{off-n-since}$). The early uploads, rehearsed on chips A and B with older code and some warm rehearsals, are
noisier (SD \val{off-sd-early}). Two checks support an offset that does not depend on the recipe. The rehearsals
since \upl{K50} span \val{since-rehearsal-span} bpb, and over that range the least-squares slope of official on
rehearsal is \val{reg-slope} (95\% CI \val{reg-lo} to \val{reg-hi}; $r=\val{reg-r}$), consistent with~1
(\cref{fig:oracle}a). The offset shows no significant drift with upload order: \val{trend-slope}$\times10^{-5}$ bpb per
upload (95\% CI \val{trend-lo} to \val{trend-hi}).

The offset was not the same on every chip (\cref{tab:offset-chips}). Chip~C, on which the calibration uploads
were rehearsed, averaged \val{chip-off-c} ($n=\val{chip-off-n-c}$); chip~H \val{chip-off-h} ($n=2$); chip~G, the
fastest final-day instance, \val{chip-off-g} ($n=3$). Within chips the SD is \val{off-sd-within}, smaller than the
pooled \val{off-sd-since}. Rounding and the run-to-run variation of the text gap explain little of the within-chip SD (\cref{app:offset-detail}). Predicting each upload
leave-one-out from the other uploads on its own chip, where there is one, gives an MAE of \val{loo-chip-mae} against
\val{loo-pooled-mae} for the pooled mean ($n=\val{loo-n}$). The offset may therefore be better described as text plus
chip, as in \cref{eq:offset}, than as a property of the text alone. With two or three uploads per new chip,
$\gamma_c$ cannot be estimated well; a step-normalised rehearsal would remove it at the source.

\paragraph{Chip or model quality?} On the final day the faster chips also rehearsed the better recipes and the
lottery winners. The uploads rehearsed on chips G and H have a mean offset \val{cq-gh-diff} above the rest, but given
the rehearsal score the G/H term is \val{cq-gh-coef} ($p=\val{cq-perm-p-adj}$; \cref{app:offset-detail}), so the
offset only \emph{may} depend on the chip.

\begin{table}[tbp]
  \centering
  \caption{\textbf{Offset by rehearsal chip since \upl{K50}} ($10^{-3}$ bpb). ``Frozen-fit error'': mean error of the
    frozen chip-C calibration (\cref{tab:offset-holdout}) on that chip's final-day uploads.
    \textsuperscript{s}\,salt chosen by a lottery; \textsuperscript{i}\,salt inherited from a lottery winner;
    \textsuperscript{n}\,rehearsal normalised to chip-C speed.}
  \label{tab:offset-chips}
  % GENERATED by paper/analysis -- offset by rehearsal chip since K50, 1e-3 bpb (oracle.py)
{\small
\begin{tabular}{@{}lp{5.4cm}rrrr@{}}
\toprule
Chip & Uploads since K50 & $n$ & Mean & SD & Frozen-fit error \\
\midrule
B & K50 & 1 & \ensuremath{+}6.86 & -- & -- \\
C & K51, K53, K54b, K56, K57, K59, K60, K63 & 8 & \ensuremath{+}6.52 & 0.39 & \ensuremath{-}0.04 \\
E & K70a\textsuperscript{n} & 1 & \ensuremath{+}6.55 & -- & \ensuremath{+}0.03 \\
H & K73s4\textsuperscript{s}, K82s4\textsuperscript{i} & 2 & \ensuremath{+}6.84 & 0.07 & \ensuremath{+}0.31 \\
G & K73s6\textsuperscript{s}, K77a\textsuperscript{i}, K82s7\textsuperscript{s} & 3 & \ensuremath{+}7.12 & 0.21 & \ensuremath{+}0.59 \\
\bottomrule
\end{tabular}
}

\end{table}

\begin{figure}[tbp]
  \centering
  \includegraphics{oracle_calibration.pdf}
  \caption{\textbf{Rehearsal and official score.} (a) Rehearsal against official for the \val{off-n-since} uploads
    since \upl{K50}; the solid line is official $=$ rehearsal \val{off-mean-since-short} with a $\pm2$\,SD band,
    the dashed line is identity. (b) The offset per upload in upload order; grey marks the early chip-A/B uploads.
    Marker shape is the rehearsal chip.}
  \label{fig:oracle}
\end{figure}

\subsection{Pseudo-prospective checks}
\label{sec:oracle-holdout}

Calibration on the uploads it describes is not a test. We therefore predict each upload only from uploads that came
before it. These analyses are computed now from the released ledger, without access to later data, so they are
leakage-free \emph{retrospective} simulations of the protocol, not archived prospective predictions.

\paragraph{Frozen calibration.} The simulator's offset model is fitted on the seven chip-C calibration uploads
\upl{K51}--\upl{K60} ($\hat\mu = \val{frozen-offset}$), the last of them scored on 29~September, before any final-day
upload. \Cref{tab:offset-holdout}
applies it, frozen, to the seven final-day uploads, which were rehearsed on four chips. The MAE is
\val{holdout-mae}~bpb (95\% bootstrap \val{holdout-mae-lo}--\val{holdout-mae-hi}), against \val{nooffset-mae} for no
offset at all, \val{priormean-mae} for the mean offset of all \val{priormean-n} earlier rehearsed uploads
(\val{priormean-offset}), \val{earlymean-mae} for the mean of the \val{earlymean-n} early chip-A/B uploads alone, and
\val{priorsince-mae} for the mean of the \val{priorsince-n} earlier uploads since \upl{K50}: the frozen fit, whose
window and calibration set were chosen with the data in view (\cref{sec:limits}), does no better than that plain
mean. The bias is \val{holdout-bias}: \val{holdout-npos}
of the 7 errors are positive (one-sided sign test $p=\val{holdout-signp}$), and the mean error is
\val{chip-ferr-c} on chip~C, \val{chip-ferr-h} on chip~H and \val{chip-ferr-g} on chip~G. This is the pattern
\cref{eq:offset} predicts for a calibration that sets $\gamma_c$ to chip C's value: the faster final-day chips
look ``hotter''.

\begin{table}[tbp]
  \centering
  \caption{\textbf{Frozen chip-C calibration} ($\hat\mu=\val{frozen-offset}$, fitted on \upl{K51}--\upl{K60}) applied to
    the final-day uploads. Error $=$ official $-$ (rehearsal $+\hat\mu$), in $10^{-3}$ bpb. Marks as in
    \cref{tab:offset-chips}. \upl{K82s4} and \upl{K82s7} were published to seven decimals
    (\val{best-official-long} and \val{salt7-official-long}).}
  \label{tab:offset-holdout}
  % GENERATED by paper/analysis -- frozen chip-C offset applied to the final-day uploads (oracle.py)
{\small
\begin{tabular}{@{}lcrrrr@{}}
\toprule
Upload & Chip & Rehearsal & Predicted & Official & Error ($10^{-3}$) \\
\midrule
K63 & C & 0.958210 & 0.96474 & 0.9647 & \ensuremath{-}0.04 \\
K70a\textsuperscript{n} & E & 0.955946 & 0.96247 & 0.9625 & \ensuremath{+}0.03 \\
K73s4\textsuperscript{s} & H & 0.955209 & 0.96174 & 0.9620 & \ensuremath{+}0.26 \\
K73s6\textsuperscript{s} & G & 0.955138 & 0.96167 & 0.9625 & \ensuremath{+}0.83 \\
K77a\textsuperscript{i} & G & 0.954729 & 0.96126 & 0.9617 & \ensuremath{+}0.44 \\
K82s4\textsuperscript{i} & H & 0.954472 & 0.96100 & 0.9614 & \ensuremath{+}0.36 \\
K82s7\textsuperscript{s} & G & 0.954379 & 0.96091 & 0.9614 & \ensuremath{+}0.49 \\
\midrule
MAE / bias & & & & & 0.35 / \ensuremath{+}0.34 \\
\bottomrule
\end{tabular}
}

\end{table}

\paragraph{Selection.} Salt lotteries select draws whose rehearsal advantage partly does not transfer
(\cref{sec:salts}), so selected uploads should miss upward. Three final-day uploads had their salt chosen by a
lottery (\upl{K73s4}, \upl{K73s6}, \upl{K82s7}; MAE \val{holdout-mae-direct}); two inherit the winning salt~4
(\upl{K77a}, \upl{K82s4}; \val{holdout-mae-inherited}); only \upl{K63} and \upl{K70a} come from no lottery
(\val{holdout-mae-none}), and one of those is normalised. Chip and selection are confounded here: two of the three
chip-G uploads are lottery winners.

\paragraph{The normalised rehearsal of \upl{K70a}.} \upl{K70a} was rehearsed on the slow chip~E and normalised to
chip-C speed with $\rho$, after the fact (\cref{sec:k70a}). With its raw rehearsal the since-\upl{K50} SD rises from
\val{sens-used-sd} to \val{sens-raw-sd} and the frozen MAE from \val{sens-used-mae} to \val{sens-raw-mae}
(\cref{tab:oracle-sens}), so the headline numbers depend on that normalisation, which is itself an application of
$\rho$.

\paragraph{Expanding window.} Predicting each upload from the mean offset of all earlier uploads since \upl{K50},
with a $t$ prediction interval and a burn-in of three, puts all \val{seq-inside} of \val{seq-n} predictions inside
their 95\% interval (\val{seq-expected} expected), with MAE \val{seq-mae} and bias \val{seq-bias}
(\cref{fig:offset-prospective}). With so few points this is weak evidence: at a burn-in of three the interval uses
$t_{2,0.975}=\val{t-two}$ and is very wide; \cref{app:oracle} gives the coverage at other levels and burn-ins.

\paragraph{The written projections.} Three projections were written before upload; \cref{sec:projections}
compares them with the official scores and explains why we treat their closeness as anecdote.

\subsection{Resolution}
\label{sec:resolution}

The oracle predicts one upload with a 95\% half-width of about \val{pi-half}~bpb ($n=\val{off-n-since}$). A
\emph{difference} between two uploads, if their errors are independent, has SD $\sqrt2\,s=\val{diff-sd}$, so the
oracle resolves official differences of about \val{diff-95}~bpb at 95\%, or about \val{diff-95-within} when both
uploads were rehearsed on the same chip. Rank order is an easier test: of the \val{conc-n} pairs of uploads since
\upl{K50}, the rehearsals ordered \val{conc-agree} the same way as the official scores (\val{conc-ties} official tie,
\val{conc-disagree} reversals), but of the \val{conc-close-n} pairs whose rehearsals differed by less than
\val{conc-close-thr}~bpb only \val{conc-close-agree} were ordered correctly. Most late levers are \val{late-lever-range}~bpb (\cref{tab:levers}). The oracle could
not have \emph{measured} them; what it did was translate chip comparisons, which do not carry the offset's noise,
onto the official scale. The levers themselves were judged on same-seed chip pairs, whose noise is discussed in
\cref{sec:seeds}.

\subsection{Traps}
\label{sec:traps}

The protocol is exact only if the rehearsal is a faithful replica of the official run. Four departures cost us
before we understood them (\cref{app:traps}): one-time compilation after the clock starts, which a warm rehearsal
does not pay (the EMA trap of \cref{sec:prewarm}); a rehearsal chip faster or slower than the official host
($\gamma_c\neq0$); a change to steps 0--19 that moves a seed onto another warm-up path, so that a measured difference
becomes a seed effect; and a time target that is a budget, not a promise. Two organiser-side anomalies are excluded
from all calibration statistics (\cref{app:traps}).

\paragraph{How we used it.} Beyond gating uploads, the oracle let us compare structural changes on chip, against a
same-chip control, without spending an upload, and run lotteries over shuffle salts and seeds; a selected draw gives
back part of its edge (\cref{sec:salts}).
